% Source: Wald GR, physical PDF p. 204
%         (printed p. 191).
% Coverage: Minkowski spacetime with a removed point on a future light cone.
% Figures/tables/footnotes: Figure 8.2; no tables or footnotes.
% Status: source-reviewed and compile-clean.
% Uncertainties: none.

\begingroup
\renewcommand{\thefigure}{8.2}
\begin{figure}[H]
  \centering
  \begin{tikzpicture}[x=.9cm,y=.9cm,>=Latex,line cap=round,font=\small]
    \draw[thick] (0,0) -- (-2.05,3.22);
    \draw[thick] (0,0) -- (2.05,3.22);
    \draw[thick] (-2.05,3.22) arc[start angle=180,end angle=360,x radius=2.05,y radius=.42];
    \draw[thick] (-2.05,3.22) arc[start angle=180,end angle=0,x radius=2.05,y radius=.42];
    \fill (0,0) circle (1.7pt) node[below] {\(p\)};
    \fill (1.72,2.52) circle (1.7pt) node[right] {\(q\)};
    \draw[fill=white] (1.26,1.77) circle (2.2pt);
    \draw[->] (1.93,1.72) .. controls (2.40,1.72) and (2.33,1.48) .. (1.42,1.70);
    \node[align=left,right] at (2.18,1.46) {(remove\\point)};
  \end{tikzpicture}
  \caption{Minkowski spacetime with a point on the future light cone of \(p\)
  removed. In this spacetime, no causal curve connects \(p\) and \(q\) so
  \(q\notin J^+(p)\) but \(q\in\overline{J^+(p)}\). Thus, \(J^+(p)\) is not
  closed.}
  \label{fig:8.2}
\end{figure}
\endgroup
